\documentclass[12pt]{article}

\usepackage[utf8]{inputenc}
\usepackage{amsmath,amssymb,mathtools}
\usepackage{booktabs}
\usepackage{float}
\usepackage{hyperref}
\usepackage[final]{microtype}

\textwidth=6.5in
\textheight=9in
\oddsidemargin=0in
\evensidemargin=0in
\topmargin=0in
\headheight=0in
\headsep=0in
\parskip=0.16in
\parindent=0in

\hypersetup{colorlinks=true, linkcolor=blue, citecolor=blue, urlcolor=blue}

\newcommand{\Tr}{\mathrm{Tr}}
\newcommand{\E}{\mathbb E}
\newcommand{\Htot}{\mathcal H_{\rm tot}}
\newcommand{\HB}{\mathcal H_{\rm B}}
\newcommand{\FR}{\mathcal F_{\rm R}}
\newcommand{\Smicro}{S_{\rm micro}}
\newcommand{\rhoB}{\rho_{\rm B}}
\newcommand{\ket}[1]{|#1\rangle}
\newcommand{\bra}[1]{\langle #1|}

\title{Schwarzschild Evaporation and Islands in a Non-Gravitational Hamiltonian Model}
\author{Sergei Slobodov\\\texttt{sergei@slobodov.com}}
\date{\today}

\begin{document}
\maketitle

\begin{abstract}
We ask which parts of Schwarzschild evaporation require gravity once the
black-hole state count is supplied as a Hamiltonian density of states.  We
construct a non-gravitational Hamiltonian with a Schwarzschild microcanonical
entropy, area-many weak emission operators, an outgoing radiation continuum,
and mixing within fixed energy shells.  The weak-coupling calculation gives
the local Hawking spectrum, the leading finite-energy correction, negative
heat capacity, and the Schwarzschild number-flux, power, and lifetime
scalings.  The area dependence of the rate comes from an inclusive sum over
boundary-accessible emission operators.  Thus the density of states fixes the
thermal factor, while the boundary emission algebra fixes the rate.  The same
Hamiltonian gives the
fine-grained information-flow structure: global purity and decoupling of the
shrinking emission map produce the Page curve and post-Page early/late
radiation correlations.  At second R\'enyi order the purity contains the two
contractions familiar from island calculations, giving the island min formula
with the microcanonical entropy of the remaining subsystem in place of the
gravitational area term.  Approximate-design and tensor-product-expander
blocks provide finite-regulator realizations of the required shell mixing.
The construction reproduces the exterior thermodynamic and entropy-flow
package without a gravitational path integral, interior geometry, or horizon
smoothness.
\end{abstract}

\section{Introduction}

We use units with $\hbar=c=k_B=1$.

This paper asks how much of Schwarzschild evaporation can be reproduced by an
ordinary quantum Hamiltonian with no dynamical geometry.  The answer we find is
that the exterior thermodynamic and entropy-flow package follows once two
black-hole inputs are supplied: the Schwarzschild microcanonical state count
and an area-sized set of weak emission operators.  Gravity enters through these
inputs.  The subsequent local thermality, rate law, Page curve, early/late
correlations, and island-formula entropy are consequences of Hamiltonian
evolution, weak emission into outgoing modes, and decoupling of the shrinking
emission map.

The target is the evaporation of an uncharged nonrotating black hole in four
asymptotically flat spacetime dimensions.  Semiclassically,
\begin{equation}
    S_{\rm BH}\propto A\propto M^2,\qquad
    T_H\propto M^{-1},
\end{equation}
and the heat capacity is negative \cite{Bekenstein1973,Hawking1975}.  The
outgoing radiation is locally thermal, up to greybody and phase-space factors
\cite{Hawking1975,Page1976}.  The luminosity scales as
\begin{equation}
    P(M)\propto M^{-2},
\end{equation}
so the lifetime scales as $\tau\propto M_0^3$.

For a unitary evaporation process, the final radiation state must encode the
initial state even though small radiation subsystems have an approximately
thermal appearance.  Page's dimension-counting argument predicts a radiation
entropy that rises and then falls \cite{Page1993Average,
Page1993Info}.  The Hayden-Preskill argument gives a sharper decoding
statement for a diary thrown into an old, rapidly scrambling black hole
\cite{HaydenPreskill2007,SekinoSusskind2008}.  The present paper uses the Page
entropy statement and early/late radiation correlations; it does not prove a
Hayden-Preskill recovery theorem.

The post-2019 island calculations add a second language for the same
fine-grained radiation entropy curve.  In those calculations, the von Neumann
entropy of the radiation is obtained by minimizing over saddles.  Before the
Page time the no-island saddle gives the Hawking entropy of the emitted
radiation.  After the Page time the island saddle gives the generalized
entropy of a quantum extremal surface, whose area term is the remaining
black-hole entropy \cite{Penington2019,
AlmheiriEngelhardtMarolfMaxfield2019,
AlmheiriHartmanMaldacenaShaghoulianTajdini2020,
AlmheiriHartmanMaldacenaShaghoulianTajdini2021}.  The Hamiltonian construction
below shows how the same min formula arises in a manifestly unitary,
non-gravitational model.  The no-island branch is the entropy of the emitted
radiation history.  The island branch is the microcanonical state count of the
remaining subsystem.  There is no gravitational path integral in the model, so
the island statement is algebraic and does not construct quantum extremal
surfaces or replica wormholes.

The exterior features recovered below are the following.
\begin{enumerate}
\item
The black-hole entropy scales as the horizon area.  For a four-dimensional
Schwarzschild black hole, $A\sim M^2$.
\item
The temperature scales as $T_H\sim M^{-1}$, so the heat capacity is negative.
\item
The outgoing radiation has an approximately thermal one-particle spectrum at
the instantaneous Hawking temperature, with the usual greybody and phase-space
factors and with finite-energy deviations from exact thermality.
\item
For four-dimensional Schwarzschild evaporation, the typical quantum energy
scales as $\omega\sim T\sim M^{-1}$, the luminosity scales as
$P\sim M^{-2}$, the number flux scales as
$\Gamma_{\rm quanta}\sim M^{-1}$, and the evaporation time scales as
$\tau\sim M_0^3$.
\item
The evaporation is quasi-stationary: over short intervals the radiation is
described by an instantaneous mass and temperature, while the macroscopic mass
changes only over many emitted quanta.
\item
The fine-grained evolution is unitary: the radiation von Neumann entropy
follows the Page curve, equivalently the island min formula for the exterior
radiation entropy, and late radiation is correlated with early radiation after
the Page time.
\end{enumerate}

The contribution of this paper is to assemble these thermodynamic, rate, and
fine-grained entropy statements in one energy-resolved Hamiltonian framework.
The Hamiltonian has a subsystem $B$ with microcanonical density of states
$\rho_{\rm B}(E)\sim \exp(cE^2)$, with the energy $E$ identified with the mass
$M$.  The leading entropy is carried by boundary-accessible degrees of
freedom.  The subsystem is weakly coupled to an outgoing radiation continuum by
an energy-conserving interaction whose local emission operators act on those
degrees of freedom.  The inclusive low-frequency emission strength is then
proportional to the number of active boundary emitters, hence to the area.
This reproduces the area scaling of the standard low-frequency absorption
cross-section by operator counting.  The same weak-emission map, after mixing
inside microcanonical shells, gives the shrinking-isometry Page curve and its
island-formula min structure.

This separation is important.  Recent work constructs operators with
black-hole-like spectral growth, $\Omega(E)\sim\exp({\rm const.}\,E^2)$, and
finds that such spectra can arise from unusual, barely localized wavefunctions
\cite{AurellMajumdar2025}.  In the present framework, the spectral growth alone
would give the microcanonical Hawking temperature.  It would not fix the
inclusive luminosity or the fine-grained Page/island structure.  Those require
the boundary-accessible emission algebra and the decoupling of the induced
shrinking channel.

In the design-based version used below, the regulated
in-shell mixing blocks are chosen from approximate-design or
tensor-product-expander constructions, so this decoupling condition follows
from standard results
\cite{HaydenPreskill2007,HaydenHorodeckiWinterYard2008,
DankertCleveEmersonLivine2009,SzehrDupuisTomamichelRenner2013,
BrownFawzi2015,HarrowLow2009}.  Known Hamiltonian-design constructions give
Hamiltonian evolutions with the required design moments under their ensemble,
locality, or time-dependence assumptions, including nearly time-independent
design Hamiltonians and time-independent random polylog-local Hamiltonians
\cite{NakataHircheKoashiWinter2017,CuiSchusterMaoHuangBrandao2025}.
Channel-scrambling diagnostics based on out-of-time-order correlator (OTOC)
decay give another Hamiltonian route to the same mutual-information statements
\cite{HosurQiRobertsYoshida2016,YoshidaKitaev2017}.

The construction is conditional in two explicit places.  First, the
Schwarzschild density of states and its boundary-accessible realization are
inputs to the Hamiltonian.  Deriving them from simpler non-gravitational
microscopic degrees of freedom is a separate problem.  Second, the
fine-grained information-flow result uses decoupling of the composite emission
map, supplied here by approximate designs or tensor-product expanders.  With
these inputs, the remaining evaporation package follows by standard
weak-coupling, microcanonical, and quantum-information arguments.  The model
has no interior geometry and does not address horizon smoothness, island
reconstruction from the radiation, complementarity, or firewall questions.

\subsubsection*{Relation to earlier work}

Several finite-dimensional models address the information-theoretic side of
this picture.  Tensor-network models, qubit-transport models, and random-circuit
models describe unitary transfer of information from a shrinking black-hole
subsystem into radiation \cite{LeutheusserVanRaamsdonk2017,Avery2013,
OsugaPage2018,PiroliSunderhaufQi2020}.  These constructions make Page-like
information flow transparent.  The Hamiltonian considered here also includes the
thermodynamic data: an energy, a density of states, emitted quanta with
definite energies, local thermal flux, negative heat capacity, and the
Schwarzschild rate law.  Related Hilbert-space reduction models make the
island-like entropy minimum explicit from changing Hilbert-space dimensions
\cite{BasuWenZhou2025}.

Analogue-gravity systems give another comparison point.  Unruh's acoustic
black holes and the later analogue-gravity literature show that
Hawking-like radiation can arise in non-gravitational media whose long-distance
excitations see an effective horizon \cite{Unruh1981,Barcelo2011}.  Those
models capture the kinematic field-theory part of Hawking radiation.  The
Hamiltonian construction below addresses the thermodynamic and
information-theoretic part by supplying the black-hole state count as a density
of states and treating evaporation as ordinary weak emission into outgoing
quantum modes.

The closest information-theoretic comparison is the tensor-network class of
Leutheusser and Van Raamsdonk, where each evaporation step is specified by an
isometry
\begin{equation}
    {\cal I}_n:{\cal H}^{BH}_n\to {\cal H}^{BH}_{n-1}\otimes{\cal H}^{rad}_n .
\end{equation}
In the present notation the corresponding map is $V_E$, used in
Section~\ref{sec:information-flow} and constructed in
Appendix~\ref{app:weak-emission-map} as the Stinespring dilation of the weak
emission channel
\begin{equation}
    V_E:\mathcal H_E\to
    \bigoplus_{\omega,\lambda,\mu}
    \mathcal H_{E-\omega}\otimes\ket{\omega,\lambda,\mu}_{\rm R}.
\end{equation}
The additional structure here is the energy-resolved Hamiltonian origin of
this map and the density of states that produces the thermodynamic evaporation
law.  Tensor-network and random-circuit models supply the Page and
early/late-correlation part of the argument
\cite{LeutheusserVanRaamsdonk2017,PiroliSunderhaufQi2020}.

Section 2 defines the model.  Section 3 derives the thermodynamic consequences.
Section 4 states the emission-isometry reduction, the Page curve, and the
island-formula contraction.  Section 5 gives the finite-dimensional
regularization needed to apply Page and decoupling estimates.  Section 6
compares the model with existing approaches, and Section~\ref{sec:discussion}
discusses the fixed-Hamiltonian problem left by the abstract mixer.  The
appendices give the
weak-emission Stinespring construction, a concrete fixed-Hamiltonian target,
the Page/replica moment calculation, and a finite-dimensional check separating
coarse thermality from Page behavior.

\section{The Model}

The model has four ingredients.  The first is a black-hole subsystem $B$,
representing the remaining object.  The second is an outgoing
radiation field $R$, representing emitted quanta that propagate away.  The
third is a mixing term that scrambles states inside fixed microcanonical
shells.  The fourth is an interaction that lowers the energy of $B$ while
creating an outgoing quantum.  The full Hamiltonian is
\begin{equation}
    H_{\rm tot}=H_{\rm B}^{(0)}+H_{\rm mix}+H_{\rm R}+H_I.
\end{equation}
The term $H_{\rm mix}$ mixes states inside a fixed microcanonical shell of
$B$.  The interaction $H_I$ transfers energy from $B$ to $R$.
The Hamiltonian is time independent.  The stochastic language used later for
emission events is the standard weak-coupling description of this Hamiltonian
in the outgoing-scattering regime.

\subsection{Black-Hole Subsystem}

The Hilbert space of the black-hole subsystem is organized by energy.  We use
continuum notation for readability, but all entropy statements refer to
microcanonical bands of width $\Delta E$ that are wide compared with the
microscopic level spacing and narrow compared with the scale on which
$\Smicro(E)$ varies:
\begin{equation}
    \HB=\int^\oplus dE\,\mathcal H_E,
    \qquad
    H_{\rm B}^{(0)}=\int dE\,E\,\Pi(E).
\end{equation}
Here $\Pi(E)$ denotes the formal spectral projector density onto energy $E$.
For a shell of width $\Delta E$, let $\Pi_{\Delta E}(E)$ denote the projector
onto states with energies in that shell.  Then
\begin{equation}
    N(E,\Delta E)=\Tr \Pi_{\Delta E}(E)\simeq \rhoB(E)\Delta E,
    \qquad
    \Smicro(E)=\log N(E,\Delta E).
\end{equation}
The Schwarzschild scaling is encoded by an area function $A(E)$ and by the
microcanonical entropy
\begin{equation}
    \Smicro(E)\simeq cE^2,
    \qquad
    A(E)\simeq c_A E^2.
\end{equation}
The constants $c$ and $c_A$ set the entropy and area units.
For the rate-law result, we use a boundary-accessible realization of this state
count.  Up to subexponential factors, the active shell contains
\begin{equation}
    N(E)\simeq \frac{\Smicro(E)}{\log q}
\end{equation}
active boundary cells of local Hilbert-space dimension $q$.  Thus
\begin{equation}
    \mathcal H_E\simeq
    \mathcal K_E\otimes
    \bigotimes_{x=1}^{N(E)}\mathcal h_x,
    \qquad
    \dim \mathcal h_x=q,
\end{equation}
where $\mathcal K_E$ denotes any remaining degeneracy that is subleading for
the operator count.  Since $\Smicro(E)\propto A(E)$, the number of active
boundary cells is area-sized:
\begin{equation}
    N(E)\propto A(E).
\end{equation}
This is the horizon-qubit bookkeeping assumption in Hamiltonian language.  It
is not a derivation of the Schwarzschild density of states from simpler
degrees of freedom.  The word ``boundary'' refers here to the degrees that are
accessible to the emission operators and that carry the leading shell entropy;
it is a bookkeeping analogue of horizon area, not a geometric bulk/boundary
factorization inside the model.
For a narrow microcanonical wavepacket, the mass variable is the expectation
value of the subsystem energy,
\begin{equation}
    M(t)=\langle H_{\rm B}^{(0)}\rangle_t.
\end{equation}
Choose a microscopic energy basis
\begin{equation}
    H_{\rm B}^{(0)}\ket a=E_a\ket a,\qquad A_a=A(E_a),
\end{equation}
with many states $\ket a$ in the same microcanonical shell.  Roman indices
$a,b$ label black-hole energy eigenstates throughout.

\subsection{Radiation Hilbert Space}

The outgoing radiation Hilbert space is a Fock space $\FR$ with Hamiltonian
\begin{equation}
    H_{\rm R}
    =
    \sum_\lambda\int_0^\infty d\omega\,
    \omega\, b_\lambda^\dagger(\omega)b_\lambda(\omega).
\end{equation}
The modes are continuum-normalized,
\begin{equation}
    [b_\lambda(\omega),b_{\lambda'}^\dagger(\omega')]
    =\delta_{\lambda\lambda'}\delta(\omega-\omega').
\end{equation}
The label $\lambda$ denotes species, angular channels, polarizations, or other
outgoing radiation labels.  Smooth density-of-modes and residual channel
factors are included below in the channel weights.  The outgoing field is
treated as an asymptotic radiation channel or large lead, so emitted quanta
propagate away from $B$ on the evaporation time scale.

\subsection{Mixing and Emission}

The mixing term acts within fixed microcanonical shells,
\begin{equation}
    H_{\rm mix}=\int dE\,H_{\rm mix}(E),
    \qquad
    H_{\rm mix}(E):\mathcal H_E\to\mathcal H_E.
\end{equation}
Equivalently, in a finite regulator one replaces the direct integral by a
direct sum over microcanonical bands and takes
$H_{\rm mix}=\oplus_E H_{\rm mix}(E)$.  This is the version used in the
dimension-counting argument below.  Exact block diagonality is an idealization
of a microcanonical description: small off-shell matrix elements may be
included provided they do not change the coarse evaporation trajectory on the
emission time scale.
For the design-based version of the model, the finite-shell mixer is
specified by a unitary block $U_E$ whose low-order moments form an approximate
design, or more generally by a tensor-product expander of the order required by
the diagnostic.  The Hamiltonian block is any self-adjoint logarithm
\begin{equation}
    H_{\rm mix}(E)=\frac{i}{t_{\rm scr}}\log U_E,
    \qquad
    e^{-iH_{\rm mix}(E)t_{\rm scr}}=U_E .
\end{equation}
Once the blocks $U_E$ are chosen, the evolution is fixed and time independent.
Approximate two-designs are enough for the second R\'enyi Page-purity and
decoupling estimates
\cite{DankertCleveEmersonLivine2009,SzehrDupuisTomamichelRenner2013};
tensor-product expanders give the corresponding higher-moment route
\cite{HarrowLow2009}.  Related Hamiltonian-design constructions are discussed
in Sec.~\ref{sec:comparison} and App.~\ref{app:fixed-hamiltonian-target}.

The weak-emission regime is
\begin{equation}
    t_{\rm emit}\ll t_{\rm evap}.
\end{equation}
Here $t_{\rm emit}$ is the mean time between emissions and $t_{\rm evap}$ is
the macroscopic evaporation time.  The information-theoretic results use the
additional mixing condition stated in Sec.~\ref{sec:isometry-flow}: the
composite emission map must act as a typical isometry on the active
microcanonical support.

The emission term uses the standard system-continuum form for weak decay into
outgoing modes, familiar from Friedrichs-Lee and Wigner-Weisskopf descriptions
of unstable quantum systems coupled to continua
\cite{Friedrichs1948,Lee1954,CohenTannoudji1992}.  In a multi-channel
Friedrichs-Lee notation, the continuum state coupled to a state $\ket a$ of
$B$ is
\begin{equation}
    \ket{b;\omega,\lambda}\equiv
    \ket b\otimes b_\lambda^\dagger(\omega)\ket{0}_{\rm R}.
\end{equation}
The interaction lowers the area associated with the black-hole subsystem and
creates an outgoing quantum.  We separate the outgoing radiation label
$\lambda$ from an inclusive emission-channel label $\mu$.  The latter
labels the active boundary cell, or equivalently an orthogonal boundary
emission operator.  The number of such labels is
$N(E)\propto A(E)$ in the boundary-accessible realization above:
\begin{align}
    H_I
    =
    \sum_{a,b,\lambda,\mu}
    \int_0^\infty d\omega\,
    \Big[
    &f^{\mu\lambda}_{ba}(\omega)
    \ket b\bra a\,
    b_\lambda^\dagger(\omega)
    +{\rm h.c.}\Big].
\end{align}
For initial states with empty outgoing radiation and for an outgoing continuum
or large lead, the weak-coupling scattering limit selects the decaying branch
on the evaporation time scale.  This is the rotating-wave, vacuum-continuum
Markov limit in which reabsorption and late-time backflow are negligible.
The form factor $f^{\mu\lambda}_{ba}(\omega)$ is supported on transitions with
$A_b<A_a$.  In the
long-time weak-coupling calculation, the usual resonant delta function selects
transitions with $E_a-E_b=\omega$.  We write the smooth part of the matrix
element as
\begin{equation}
    f^{\mu\lambda}_{ba}(\omega)
    =
    g\,u^{\mu\lambda}_{ba}\,
    \omega^{p/2}\,
    \chi(A_b,A_a;\omega).
\end{equation}
Here $g$ is the weak coupling constant.  The dimensionless matrix element
$u^{\mu\lambda}_{ba}$ distributes amplitude among available final states, and $\chi$
restricts the allowed area loss.  The factor
$\omega^p$ is the effective low-energy phase-space weight; for massless
emission in three spatial dimensions the leading behavior is $p=2$.  The
area factor enters after summing over the active boundary emitters.  For
four-dimensional Schwarzschild emission the corresponding low-frequency
absorption cross-section is proportional to the horizon area,
$\sigma_{\rm abs}(\omega\to0)=A_H$ for minimally coupled massless scalars
\cite{Page1976,DasGibbonsMathur1997}.  The present Hamiltonian reproduces the
same area scaling through an inclusive sum over area-many weak boundary
emitters.  The corresponding channel count is
\begin{equation}
    N_A(E)\equiv \#\{\mu\}=N(E)\propto A(E).
\end{equation}
The channel sum is an inclusive sum over independent weak boundary emitters.
The channels are orthogonal in the outgoing continuum, independently phased,
or shell-averaged, so their probabilities add in the golden-rule rate.
Coherent addition of identical amplitudes would give a different scaling and is
not the regime considered here.
Residual greybody and channel dependence is placed in the smooth functions
$C_\lambda(E,\omega)$ below.  The Boltzmann factor comes from the
density-of-states ratio.  Other values of $p$ or different scaling of
$N(E)$ give models with the same microcanonical thermal factor and a different
total power law.

The same statement can be written as an inclusive spectral function of the
boundary emission operators:
\begin{equation}
    {\cal A}_\lambda(E,\omega)
    =
    \frac{1}{D_E}
    \sum_\mu
    \Tr\!\left[
    \Pi_E O_{\lambda\mu}^\dagger(\omega)
    \Pi_{E-\omega}
    O_{\lambda\mu}(\omega)\Pi_E
    \right].
\end{equation}
Here $O_{\lambda\mu}(\omega)$ is the part of the interaction operator whose
matrix elements are $u^{\mu\lambda}_{ba}$ between the shells $E$ and
$E-\omega$.
If the active boundary emitters have comparable shell-averaged strength and
their cross terms vanish in the inclusive rate, then
\begin{equation}
    {\cal A}_\lambda(E,\omega)
    \simeq
    N(E) C_\lambda(E,\omega)
    \frac{D_{E-\omega}}{D_E}.
\label{eq:boundary-spectral-function}
\end{equation}
Thus the area factor is an operator-counting consequence of boundary-accessible
entropy and local weak emission.  This area factor is used in the rate law; the
Page and island calculations below use the same emission labels as radiation
records, but their leading entropy competition is controlled by Hilbert-space
dimensions and decoupling.

For a rapidly mixing subsystem, the emission matrix elements are taken to have
the standard ETH/random-matrix scaling in the energy basis
\cite{Deutsch1991,Srednicki1994}.  In the usual ETH notation the off-diagonal
matrix element carries a factor exponential in one half of the entropy of the
relevant shell.  Here the transition begins in shell $E$.  For $\omega\ll E$,
and for $a$ in shell $E$ and $b$ in shell $E-\omega$, we write
\begin{equation}
    u^{\mu\lambda}_{ba}=D_E^{-1/2} r^{\mu\lambda}_{ba},
    \qquad D_E=\dim \mathcal H_E,
\end{equation}
where the dimensionless variables $r^{\mu\lambda}_{ba}$ have zero mean and shell-averaged
variance
\begin{equation}
    \overline{|r^{\mu\lambda}_{ba}|^2}=C_\lambda(E,\omega).
\end{equation}
The overline denotes the shell or ensemble average within the indicated
microcanonical bands.
The rate calculation uses only this shell-averaged normalization.  The
information-theoretic statements below require the composed emission channel to
decouple the reference from the complementary subsystem.  This is the
channel-scrambling input supplied by the shell mixer.
The smooth function $C_\lambda(E,\omega)$ includes residual channel-dependent
greybody and form-factor dependence not already contained in the phase-space
factor or in the area cross-section.  This
normalization gives
\begin{equation}
    \sum_{b\in E-\omega}|u^{\mu\lambda}_{ba}|^2
    =
    C_\lambda(E,\omega)
    \frac{D_{E-\omega}}{D_E}
    \left[1+O(D_{E-\omega}^{-1/2})\right].
\label{eq:inclusive-u}
\end{equation}
Thus the inclusive rate from a typical initial microstate depends on shell
degeneracies and is insensitive to special basis choices inside the shell.  The
sum over $\mu$ gives the area-sized absorption factor, while the smooth
function $C_\lambda(E,\omega)$ contains residual channel dependence.

\subsection{Initial States and Observables}

The initial states are pure states of $B$ in a narrow microcanonical band, with
empty outgoing radiation:
\begin{equation}
    \ket{\Psi(0)}\in \mathcal H_E\otimes\ket{0}_{\rm R}.
\end{equation}
We consider typical states in the microcanonical shell and states obtained by
acting with a low-complexity perturbation on such states.

The basic observables are:
\begin{align}
    M(t)&=\langle H_{\rm B}^{(0)}\rangle_t,\\
    A(t)&=\langle A(H_{\rm B}^{(0)})\rangle_t,\\
    E_{\rm rad}(t)&=\langle H_{\rm R}\rangle_t,\\
    P(t)&=\frac{d}{dt}E_{\rm rad}(t),
\end{align}
together with the one-particle radiation spectrum.  The information-theoretic
observables are radiation entanglement entropies and mutual information, for
example
\begin{align}
    S_{\rm vN}(R)&=-\Tr\rho_R\log\rho_R,\\
    I(R_{\rm early}:R_{\rm late})
    &=S(R_{\rm early})+S(R_{\rm late})-S(R_{\rm early}R_{\rm late}).
\end{align}
The entropy $\Smicro$ is the microcanonical state-counting entropy of $B$.
The entropy $S_{\rm vN}$ is the von Neumann entropy of a radiation subsystem.

\section{Thermodynamic Properties}

The density of states of $B$ fixes the microcanonical temperature.  For
\begin{equation}
    \Smicro(M)=cM^2,
\end{equation}
the inverse temperature is
\begin{equation}
    \beta(M)=\frac{d\Smicro}{dM}=2cM.
\end{equation}
Thus
\begin{equation}
    T(M)=\beta(M)^{-1}\propto M^{-1},
\end{equation}
and
\begin{equation}
    C=\frac{dM}{dT}<0.
\end{equation}

The local emission spectrum follows from Fermi's golden rule and the
microcanonical degeneracy ratio.  For an emitted quantum of energy
$\omega\ll M$,
\begin{equation}
    \frac{\rhoB(M-\omega)}
         {\rhoB(M)}
    =
    \exp[\Smicro(M-\omega)-\Smicro(M)]
    =
    \exp[-\beta(M)\omega+c\omega^2]
    \simeq e^{-\beta(M)\omega}.
\end{equation}
The final approximation is the local thermal limit.  The retained
$c\omega^2$ term is the leading finite-energy backreaction correction, the
same entropy-difference correction emphasized in the tunneling description of
Parikh and Wilczek \cite{ParikhWilczek2000}.  For typical Hawking quanta
$\omega\sim T\sim M^{-1}$, this correction is subleading at large $M$ but is
the leading non-thermal deviation from an exactly canonical spectrum.
In the weak-coupling continuum regime, Fermi's golden rule gives a transition
rate equal to the squared matrix element times the density of final states
\cite{Friedrichs1948,CohenTannoudji1992}.  With the shell-averaged
matrix-element normalization in Eq.~\eqref{eq:inclusive-u}, and after summing
over the active boundary emitters, this gives
\begin{equation}
    \frac{d\Gamma}{d\omega}
    \propto
    N(M)\,\omega^p\,
    \exp[\Smicro(M-\omega)-\Smicro(M)].
\end{equation}
For typical Hawking quanta, this becomes
\begin{equation}
    \frac{d\Gamma}{d\omega}
    \propto N(M)\omega^p e^{-\beta(M)\omega}.
\end{equation}

The particle emission rate and emitted power scale as
\begin{align}
    \Gamma_{\rm quanta}(M)
    &\propto
    N(M)
    \int_0^\infty d\omega\,\omega^p e^{-\beta(M)\omega}
    \propto N(M)\beta(M)^{-(p+1)},\\
    P(M)
    &\propto
    N(M)
    \int_0^\infty d\omega\,\omega^{p+1}e^{-\beta(M)\omega}
    \propto N(M)\beta(M)^{-(p+2)}.
\end{align}
With $N(M)\propto A(M)\propto M^2$, $\beta(M)\propto M$, and $p=2$,
\begin{equation}
    \Gamma_{\rm quanta}(M)\propto M^{-1},
    \qquad
    P(M)\propto M^{-2}.
\end{equation}
The local thermal factor and its leading finite-energy correction follow from
$\rhoB(M-\omega)/\rhoB(M)$ alone.  The total Hawking power law also uses the
area-sized boundary operator count.  If only order-one boundary emitters
couple efficiently, the three-spatial-dimensional phase-space scaling becomes
$P\propto M^{-4}$.
The mass therefore obeys
\begin{equation}
    \frac{dM}{dt}\propto -M^{-2},
\end{equation}
and the lifetime scales as
\begin{equation}
    \tau\propto M_0^3.
\end{equation}
The quasi-stationary expansion fails near the endpoint, where the emitted
quantum energy is no longer small compared with the remaining mass.  The
calculation applies in the large-mass evaporation regime, as in the standard
semiclassical treatment.

Thus the thermodynamic part of Schwarzschild evaporation follows from the
boundary-accessible density of states and local weak emission.  The
density-of-states ratio gives local thermality and the leading finite-energy
correction.  The boundary operator count fixes the four-dimensional power law.

\section{Isometry Reduction and Information Flow}
\label{sec:isometry-flow}
\label{sec:information-flow}

The thermodynamic calculation uses inclusive transition rates.  Fine-grained
information flow uses the corresponding maps between microcanonical Hilbert
spaces.  In the weak-coupling Markov regime, conditioning on one emitted
quantum gives an emission isometry
\begin{equation}
    V_E:\mathcal H_E\longrightarrow
    \bigoplus_{\omega,\lambda,\mu}
    \mathcal H_{E-\omega}\otimes \ket{\omega,\lambda,\mu}_{\rm R}.
\label{eq:emission-isometry-main}
\end{equation}
The labels record the emitted energy, outgoing channel, and boundary emission
operator.  The branch weights are the golden-rule weights derived above,
including the density-of-states ratio
$D_{E-\omega}/D_E$ and the area-sized sum over active boundary emitters.  The
standard quantum-jump construction of Eq.~\eqref{eq:emission-isometry-main}
and the waiting-time record are given in Appendix~\ref{app:weak-emission-map}.

For a narrow coarse-grained evaporation trajectory
\begin{equation}
    E_0\to E_1\to \cdots \to E_m=E,
\end{equation}
the one-step maps compose to
\begin{equation}
    V_{E_0\to E}:
    \mathcal H_{E_0}\longrightarrow
    \mathcal H_E\otimes \mathcal H_{R(E_0\to E)} .
\label{eq:composite-isometry-main}
\end{equation}
The radiation Hilbert space in Eq.~\eqref{eq:composite-isometry-main} carries
the emitted energies, channel labels, boundary labels, and time-bin labels
along the trajectory.  Energy conservation fixes the remaining shell, while
the density-of-states ratio fixes the thermodynamic branching weights.

The fine-grained assumption is a typical-encoding condition for the composite
map on the active microcanonical support.  In the regulated model this
condition is supplied by the approximate-design or tensor-product-expander
mixing blocks described in Section~2.  Page's theorem and standard decoupling
estimates then apply to the factorization
$\mathcal H_E\otimes \mathcal H_{R(E_0\to E)}$
\cite{Page1993Average,Page1993Info,HaydenHorodeckiWinterYard2008}.
Operationally, this means the following.  Let $X_E$ be the smaller factor in
the bipartition
$\mathcal H_E\otimes\mathcal H_{R(E_0\to E)}$.  For the code subspace under
consideration, the reduced state on $X_E$ is close to the corresponding
Page-typical state: maximally mixed in the equiprobable-history idealization,
or Hawking-weighted on the radiation branch when the golden-rule history
weights are kept.  Denote the trace-distance error by $\epsilon_{\rm Page}(E)$.
Then the entropy errors are bounded by the usual continuity estimates and are
absorbed in the $O(1)$ term below in the large-entropy regime.  Equivalently,
for code states purified by a reference, the complementary subsystem is
decoupled from that reference, up to the same error.

Let
\begin{equation}
    \Delta S_{\rm rad}(E)=\Smicro(E_0)-\Smicro(E),
    \qquad
    D_{\rm B}(E)=e^{\Smicro(E)}.
\end{equation}
Here $\Delta S_{\rm rad}$ is the coarse-grained thermodynamic entropy carried
by the emitted Hawking history, and $D_{\rm B}$ is the remaining core state
count.  Under the typical-encoding condition, the exterior radiation entropy is
\begin{equation}
    S_{\rm vN}(R;E)
    =
    \min\{\Delta S_{\rm rad}(E),\Smicro(E)\}
    +O(1).
\label{eq:pagecurve}
\end{equation}
The post-Page decreasing branch also follows directly from global purity:
$S_{\rm vN}(R)=S_{\rm vN}(B)\le \Smicro(E)$ once the remaining subsystem is the
smaller factor.  The typical-encoding condition is what saturates the smaller
factor on both sides of the Page transition.

The same statement has a replica form.  At second R\'enyi order, for
equiprobable radiation histories with effective dimension
$D_{\rm rad}=e^{\Delta S_{\rm rad}}$,
\begin{equation}
    \Tr\rho_R^2
    \simeq
    e^{-\Delta S_{\rm rad}(E)}
    +
    e^{-\Smicro(E)} .
\label{eq:renyi-two-island}
\end{equation}
For Hawking-weighted histories, labelled by probabilities $p_\alpha$, the
first term is replaced by $\Tr p^2$:
\begin{equation}
    \Tr\rho_R^2
    \simeq
    \Tr p^2+e^{-\Smicro(E)} .
\label{eq:weighted-renyi-two-main}
\end{equation}
Appendix~\ref{app:page-replica} gives the Page and permutation-moment
calculation behind Eqs.~\eqref{eq:renyi-two-island}
and~\eqref{eq:weighted-renyi-two-main}.

Equation~\eqref{eq:pagecurve} is the island formula in microcanonical
variables.  The first entry is the no-island Hawking entropy of the emitted
radiation history.  The second entry is the island branch, written here as the
logarithm of the remaining microcanonical Hilbert-space dimension.  In the
gravitational quantum-extremal-surface formula this second term is the
generalized entropy, with area contribution $A/4$ equal to the black-hole
entropy at the later mass \cite{Penington2019,
AlmheiriEngelhardtMarolfMaxfield2019,
AlmheiriHartmanMaldacenaShaghoulianTajdini2021}.  Here it is
$\log D_{\rm B}(E)$.  The second R\'enyi contraction is the corresponding
non-geometric replica contraction: replicas connect through the smaller
remaining core factor, with no spacetime saddle.

The correspondence is therefore:
\begin{center}
\begin{tabular}{ll}
\toprule
gravitational island calculation & Hamiltonian model \\
\midrule
radiation region & emitted radiation record \\
no-island saddle & Hawking-history contraction \\
island saddle & remaining-core contraction \\
area term $A/4G$ & $\log D_{\rm B}(E)=\Smicro(E)$ \\
QES minimization & minimization over the remaining shell \\
replica wormhole saddle & non-geometric replica contraction \\
\bottomrule
\end{tabular}
\end{center}
This table is a dictionary for the entropy calculation only.  The Hamiltonian
model has no spacetime region bounded by a quantum extremal surface and no
gravitational path integral.

Combining the finite-dimensional Page estimate with the regulated
weak-emission construction gives the schematic accuracy statement
\begin{equation}
    S_{\rm vN}(R;E)
    =
    \min\{\Delta S_{\rm rad}(E),\Smicro(E)\}
    +O(1)
    +O\!\left(\epsilon_{\rm Page}(E)\log D_{\rm min}(E)\right)
    +\delta_{\rm shell}+\delta_{\rm gb}+\delta_{\rm end}.
\end{equation}
Here $D_{\rm min}(E)=\min\{D_{\rm B}(E),e^{\Delta S_{\rm rad}(E)}\}$.
The first correction is the finite-dimensional Page correction.  The second is
the error from imperfect decoupling of the composite emission map.  The terms
$\delta_{\rm shell}$, $\delta_{\rm gb}$, and $\delta_{\rm end}$ denote,
respectively, finite shell-width errors, greybody/history-weighting errors,
and the breakdown of the quasi-stationary large-mass expansion near the
endpoint.

The Page time is determined by
\begin{equation}
    \Smicro(E_{\rm Page})
    =
    \frac12\Smicro(E_0),
\end{equation}
up to corrections set by shell width and greybody-weighted branching.  After
this time, a late radiation block is locally close to thermal while its
purifier is mostly the early radiation.  Thus
\begin{equation}
    I(R_{\rm early}:R_{\rm late})
    =
    S(R_{\rm early})+S(R_{\rm late})-S(R_{\rm early}R_{\rm late})
\end{equation}
becomes nonzero after the Page time, because
$S(R_{\rm early}R_{\rm late})=S(B)$ by global purity and $S(B)$ is bounded by
the smaller remaining Hilbert space.  This is a Page-level correlation
statement.  A Hayden-Preskill recovery theorem would require an additional
diary system and a sharper dynamical scrambling estimate.

The fine-grained result therefore has two layers.  The Hamiltonian and
shell-averaged matrix elements give the density-of-states ratio in the
inclusive rates.  Global purity and the shrinking remaining Hilbert space give
the post-Page entropy bound.  Saturation of the rising branch, saturation of
the post-Page bound, and early/late correlations use decoupling of the
composite emission map on the active microcanonical support.  The weak-coupling
and quasi-stationary derivation applies while the emitted quantum energy is
small compared with the remaining energy.

\section{Finite-Dimensional Regulator}

This section makes the finite-dimensional part of the argument explicit.  The
Hamiltonian calculation uses continuous energy and radiation labels, while
Page and decoupling estimates apply to finite Hilbert spaces.  The regulator
below replaces the continuous labels by finite bands and wave packets while
preserving the weak-coupling emission branch.  Appendix~\ref{app:finite-check}
gives a small finite-dimensional check of the Page-theorem step and of the
distinction between coarse energy-bin thermality and fine-grained radiation
entropy.

The continuum labels are replaced by finite labels in two steps.  Black-hole
energies are grouped into microcanonical bands.  Outgoing radiation modes are
grouped into wave packets with finite frequency, channel, and time-bin labels.
After this replacement, each emission branch is a finite map from one
microcanonical Hilbert space to a smaller microcanonical Hilbert space tensored
with a radiation label.  The continuum formulas are recovered by refining the
bands and wave packets.

Choose energy bands $I_i=[E_i,E_i+\Delta E]$, with $\Delta E$ large compared
with the microscopic level spacing and small on the scale over which
$\Smicro(E)$ changes.  The black-hole Hilbert space becomes
\begin{equation}
    \HB^{\rm reg}
    =
    \bigoplus_i \mathcal H_i,
    \qquad
    \dim\mathcal H_i
    =
    D_i
    \simeq \rhoB(E_i)\Delta E .
\end{equation}
The radiation continuum is regulated by wave-packet modes labelled by a
frequency bin $\alpha$, an outgoing radiation label $\lambda$, an inclusive
emission-channel label $\mu$, and an outgoing time bin $\tau$,
\begin{equation}
    \ket{\omega,\lambda,\mu}_{\rm R}
    \quad\longrightarrow\quad
    \ket{\alpha,\lambda,\mu,\tau}_{\rm R}.
\end{equation}
The bin dimensions and weights are chosen so that sums over
$\alpha,\lambda,\mu,\tau$ converge to the continuum integrals in the
weak-coupling rate formula.

With this regulator, the emission branch of one weak-coupling time step is a
finite Stinespring map
\begin{equation}
    V_i^{\rm reg}:
    \mathcal H_i
    \longrightarrow
    \bigoplus_{\alpha,\lambda,\mu}
    \mathcal H_{i-\alpha}\otimes
    \ket{\alpha,\lambda,\mu,\tau}_{\rm R},
\end{equation}
where $\omega_\alpha$ is a representative frequency in bin $\alpha$, and
$i-\alpha$ denotes the band containing $E_i-\omega_\alpha$.  Its branch
weights are the discretized versions of $w_E(\omega,\lambda,\mu)$, including
the density-of-states ratio and the area-sized emission-channel count.
Refining the energy bands, frequency bins, and time bins recovers the
continuum rate calculation, while keeping finite-dimensional spaces at every
intermediate stage.

The regulated map is the finite-dimensional object to which Page's theorem is
applied.  The regulator preserves the weak-coupling emission branch and replaces
the continuous energy and radiation labels by finite bands and wave packets.
This is the only role of the regulator in the main argument.  It supplies the
finite-dimensional setting in which Page and decoupling statements are
applied to the Hamiltonian emission map.

\section{Comparison with Existing Models}
\label{sec:comparison}

Tensor-network, qubit-transport, and random-circuit models give explicitly
unitary descriptions of information transfer from a shrinking black-hole
subsystem into radiation
\cite{LeutheusserVanRaamsdonk2017,Avery2013,OsugaPage2018,
PiroliSunderhaufQi2020,HottaNambuYamaguchi2018}.  Their basic dynamical object
is an isometry between black-hole and radiation Hilbert spaces.  Here that
isometry is the weak-coupling emission map of an energy-resolved
Hamiltonian.

Recent non-gravitational models show that Page-like entropy dynamics can occur
in ordinary quantum systems, including open-system models, exactly solvable
free-fermion dynamics, and interacting spin chains
\cite{Glatthard2024,Glatthard2025,Kehrein2024,RayDharKulkarni2025}.  These
models separate Page-like entropy dynamics from strictly gravitational
mechanisms.  The present construction ties the same information-flow logic to
a Schwarzschild density of states, an energy-conserving emission channel,
and the Schwarzschild rate scalings.

Quantum-optical models of unitary black-hole evaporation track energy transfer
and approximate thermality in explicitly unitary models
\cite{Alsing2025,Alsing2026}.  Related recent protocol models use
coherently branched evaporation histories or time-dependent interior-radiation
pair dynamics to maintain microscopic unitarity
\cite{AkilGiannelliModestoDahlstenChiribellaBrukner2025,Abutaleb2026}.  The
thermal factor here comes from the microcanonical density-of-states ratio of
the black-hole subsystem.
Local spin-chain models with a dynamically changing system/environment split
also reproduce Page-curve behavior and show how much of that behavior can be
kinematic \cite{JonesAltaieVarcoe2026}.  The Hamiltonian here connects the
shrinking state count to emitted energy and rate data as well as to the
radiation entropy curve.
Matrix-theory black-hole models provide a more microscopic route to clump
evaporation and decreasing active degrees of freedom, with a different
holographic motivation
\cite{BanksFischlerKlebanovSusskind1998,BerkowitzHanadaMaltz2016}.

The Hamiltonian class uses four ingredients: boundary-accessible Schwarzschild
state counting, an outgoing continuum, local weak emission from the active
boundary cells, and design-based decoupling of the composite emission map.
The area-sized inclusive strength follows from the boundary operator count.
This combination ties Page's isometry logic to the Schwarzschild thermodynamic
scalings.  The radiation entropy also has the island-formula structure, and
its second R\'enyi purity has the corresponding two-term replica competition,
with the gravitational area term replaced by the microcanonical entropy of the
remaining subsystem.

\section{Discussion}
\label{sec:discussion}

The model separates the exterior evaporation package into a small number of
Hamiltonian ingredients.  The density of states fixes the thermodynamic
behavior: temperature, negative heat capacity, the local thermal factor, and
the leading finite-energy correction.  Boundary accessibility and local weak
coupling turn that state count into an area-sized inclusive emission strength,
which fixes the Schwarzschild number flux, power, and lifetime scalings.  The
fine-grained entropy statements use global purity, the shrinking remaining
Hilbert space, and decoupling of the composite emission map.

The island comparison is algebraic.  The no-island branch is the entropy of
the emitted Hawking history.  The island branch is the logarithm of the
remaining microcanonical state count.  In gravitational calculations that
second term is the generalized entropy of a quantum extremal surface; in this
Hamiltonian model it is $\log D_{\rm B}(E)$.  The second R\'enyi calculation
has the same two replica contractions: the post-Page contraction connects
replicas through the smaller remaining core factor.

The construction is conditional in two places.  First, the Schwarzschild
density of states is an input to the Hamiltonian.  A more microscopic
non-gravitational model would derive this boundary-accessible state count from
simpler degrees of freedom.  Second, the main text uses approximate-design or
tensor-product-expander mixing to supply the decoupling
needed for the Page curve and early/late correlations.  Known
Hamiltonian-design constructions provide related Hamiltonian-generated
versions under their own ensemble, locality, or timing assumptions
\cite{NakataHircheKoashiWinter2017,CuiSchusterMaoHuangBrandao2025,
ZhouZhouSonner2026}.

The natural fixed-Hamiltonian problem is more specific than asking for a
generic scrambler.  The weak-emission block must decouple the code while it
maps one shell to a smaller shell.  Appendix~\ref{app:fixed-hamiltonian-target}
spells out one concrete Cayley-expander target.  Vertex transitivity makes the
boundary channels equivalent by symmetry; the remaining dynamical requirement
is sectorwise second-moment decoupling for the open weak-emission channel,
with exact symmetry charges tracked or recorded in the radiation.  This is the
main route by which the abstract mixer could be replaced by a simple
deterministic sparse Hamiltonian.

The scope is exterior entropy flow.  The model has no interior geometry and
does not address horizon smoothness, reconstruction of an island from the
radiation, complementarity, or firewall questions.  It also does not control
the endpoint of evaporation, where the emitted quantum energy is no longer
small compared with the remaining energy and the weak-coupling quasi-stationary
expansion breaks down.  Within the stated Hamiltonian class and away from the
endpoint, local thermality, finite-energy backreaction corrections, negative
heat capacity, the Schwarzschild rate laws, the island-formula Page curve, and
post-Page early/late radiation correlations all follow from non-gravitational
Hamiltonian evolution, weak-coupling emission, microcanonical state counting,
global purity, and design-based decoupling.

\appendix

\section{Weak-Emission Stinespring Map}
\label{app:weak-emission-map}

This appendix records the weak-coupling construction behind
Eq.~\eqref{eq:emission-isometry-main}.  The radiation continuum is assumed to
have a short correlation time compared with $t_{\rm emit}$, and the coupling is
weak enough that successive emissions are well separated
\cite{CohenTannoudji1992,BreuerPetruccione2002}.  The Markov approximation then
gives an outgoing channel, while the exponentially dense bands of $B$ supply
the initial and final system subspaces.

For an initial state typical in shell $E$, Fermi's golden rule gives, for a
fixed inclusive boundary channel $\mu$,
\begin{equation}
    \frac{d\Gamma_{E\to E-\omega,\lambda\mu}}{d\omega}
    =
    2\pi g^2 \omega^p
    C_\lambda(E,\omega)
    \frac{D_{E-\omega}}{D_E}.
\end{equation}
After summing over the $N(E)$ active boundary emitters,
\begin{equation}
    \sum_\mu\frac{d\Gamma_{E\to E-\omega,\lambda\mu}}{d\omega}
    =
    2\pi g^2 N(E)\omega^p C_\lambda(E,\omega)
    \exp[\Smicro(E-\omega)-\Smicro(E)],
\end{equation}
up to the overall normalization of the boundary operator count.

For a small interval $dt\ll t_{\rm emit}$, the emission Kraus operators are
\begin{equation}
    K_{E;\omega\lambda\mu}:\mathcal H_E
    \longrightarrow \mathcal H_{E-\omega},
    \qquad
    \bra bK_{E;\omega\lambda\mu}\ket a
    =
    \sqrt{dt}\,\bra bL_{E;\omega\lambda\mu}\ket a,
\end{equation}
with
\begin{equation}
    L_{E;\omega\lambda\mu}
    =
    \sqrt{2\pi}\,g\,\omega^{p/2}
    \sum_{a\in E,\ b\in E-\omega}
    u^{\mu\lambda}_{ba}\ket b\bra a .
\end{equation}
The no-emission Kraus operator is
\begin{equation}
    K_{E;0}
    =
    I_{\mathcal H_E}
    -
    \frac{dt}{2}
    \sum_{\lambda,\mu}\int d\omega\,
    L_{E;\omega\lambda\mu}^\dagger L_{E;\omega\lambda\mu}
    +O(dt^2),
\end{equation}
so
\begin{equation}
    K_{E;0}^\dagger K_{E;0}
    +
    \sum_{\lambda,\mu}\int d\omega\,
    K_{E;\omega\lambda\mu}^\dagger K_{E;\omega\lambda\mu}
    =
    I_{\mathcal H_E}+O(dt^2).
\end{equation}

Define the total decay operator and its shell-averaged rate by
\begin{equation}
    J_E
    =
    \sum_{\lambda,\mu}\int d\omega\,
    L_{E;\omega\lambda\mu}^\dagger L_{E;\omega\lambda\mu},
    \qquad
    \Gamma_E=\frac{1}{D_E}\Tr_{\mathcal H_E}J_E .
\end{equation}
Conditioning on an emission gives normalized jumps
\begin{equation}
    \widehat K_{E;\omega\lambda\mu}
    =
    \Gamma_E^{-1/2}L_{E;\omega\lambda\mu}.
\end{equation}
If $J_E=\Gamma_EI_{\mathcal H_E}$, then
\begin{equation}
    \sum_{\lambda,\mu}\int d\omega\,
    \widehat K_{E;\omega\lambda\mu}^\dagger
    \widehat K_{E;\omega\lambda\mu}
    =
    I_{\mathcal H_E}.
\end{equation}
More generally, write $J_E=\Gamma_E(I+F_E)$.  The flat-rate map is the leading
conditioned channel, and $F_E$ measures branch-weight and timing leakage.
With the ETH/random-matrix normalization used in the main text, $F_E$ is small
on typical microcanonical states, up to shell fluctuations and Markov errors.

The conditioned Stinespring map is
\begin{equation}
    V_E=
    \sum_{\lambda,\mu}\int d\omega\,
    \widehat K_{E;\omega\lambda\mu}
    \otimes\ket{\omega,\lambda,\mu}_{\rm R}.
\end{equation}
Including the waiting-time record, the flat-rate one-emission isometry is
\begin{equation}
    W_E\ket\psi
    =
    \sum_{\lambda,\mu}\int d\omega\int_0^\infty dt\,
    \sqrt{\Gamma_E}\,e^{-\Gamma_Et/2}
    \widehat K_{E;\omega\lambda\mu}
    e^{-iH_{\rm mix}(E)t}\ket\psi
    \otimes\ket{\omega,\lambda,\mu,t}_{\rm R}.
\end{equation}
In the eigenbasis $H_{\rm mix}(E)\ket\nu=\theta_\nu\ket\nu$, tracing the
waiting-time label gives
\begin{equation}
    (\mathcal D_{\Gamma_E}\sigma)_{\nu\nu'}
    =
    \frac{\Gamma_E}{\Gamma_E+i(\theta_\nu-\theta_{\nu'})}
    \sigma_{\nu\nu'} .
\end{equation}
This Lorentzian dephasing is the temporal average supplied by the physical
emission process.  Coherences with
$|\theta_\nu-\theta_{\nu'}|\lesssim\Gamma_E$ survive with order-one weight and
belong to the remaining deterministic-channel problem discussed in
Section~\ref{sec:discussion}.

\section{A Fixed-Hamiltonian Target}
\label{app:fixed-hamiltonian-target}

This appendix states one concrete version of the fixed sparse Hamiltonian
problem mentioned in Section~\ref{sec:discussion}.  It is not used in the main
derivation, which relies on design-based mixing blocks.  Its role is to make
the remaining deterministic problem precise.

Let the active boundary cells be the vertices of a Cayley expander
$G_N={\rm Cay}({\cal G}_N,S_N)$, with homogeneous spin interactions on the
edges.  A representative shell Hamiltonian is
\begin{equation}
    K_N=
    \sum_{g\in{\cal G}_N}\sum_{s\in S_N^+}
    \left(
    J_s^xX_gX_{gs}+J_s^yY_gY_{gs}+J_s^zZ_gZ_{gs}
    \right)
    +\sum_{g\in{\cal G}_N}(h_xX_g+h_zZ_g),
\end{equation}
with no resampling during evaporation.  Cayley graphs are vertex-transitive,
and explicit Cayley/Ramanujan expander families are standard
\cite{LubotzkyPhillipsSarnak1988,Harrow2008}.  If the shell projector commutes
with the left group action and the boundary emission operators are related by
that action,
\begin{equation}
    O_{ag,\lambda}=U_aO_{g,\lambda}U_a^\dagger,\qquad
    [U_a,K_N]=[U_a,\Pi_E]=0,
\end{equation}
then the shell-averaged emission weights are equal for all boundary labels:
\begin{equation}
    {\cal A}_{g,\lambda}(E,\omega)
    =
    {\cal A}_{g',\lambda}(E,\omega).
\end{equation}
Thus the area factor follows from an exact symmetry of the boundary algebra,
while the common spectral envelope is the usual ETH input for a chaotic local
operator \cite{Srednicki1994,DAlessioKafriPolkovnikovRigol2016,
DymarskyLashkariLiu2018}.  Exact symmetry charges should be treated within
fixed sectors or recorded in the radiation.

The Page-purity and second R\'enyi island diagnostics require second-moment
decoupling of the composed evaporation map.  A sectorwise target is
\begin{equation}
    I_2(Q:B_{j+1})
    \le
    \kappa_j I_2(Q:B_j)+\varepsilon_{{\rm mix},j},
    \qquad
    \kappa_j<1,
\end{equation}
where $I_2$ is the second R\'enyi mutual information between a reference $Q$
and the remaining subsystem.  The contraction comes from the dimension drop
and from in-shell mixing that spreads reference information over the boundary
algebra.

For a fixed Hamiltonian, closed evolution by $K_N$ alone preserves
Hilbert-Schmidt distances.  The contraction relevant to Page decoupling
belongs to the weak-emission block generated by $K_N$ and $H_I$, after the
emitted mode is recorded as radiation or traced from the remaining core.  In
the flat-rate Markov limit, the one-emission block is controlled by the record
algebra
\begin{equation}
    \widehat K_{m'}^\dagger\widehat K_m,
    \qquad
    m=(\omega,\lambda,\mu),
\end{equation}
and by the emission-image overlap kernel
\begin{equation}
    G_{\nu\nu'}
    =
    \Tr\!\left[
    \rho'(\nu)\rho'(\nu')
    \right],
    \qquad
    \rho'(\nu)
    =
    \sum_m\widehat K_m\ket\nu\bra\nu\widehat K_m^\dagger .
\end{equation}
A sufficient deterministic condition is that, within each exact symmetry
sector, the diagonal matrix elements of the record algebra are smooth in the
$K_N$ eigenbasis, the kernel $G_{\nu\nu'}$ is close to its microcanonical value
for code-relevant eigenstates, and the waiting-time dephasing leaves negligible
coherent weight on near-resonant pairs.

This is a shrinking-channel decoupling estimate, not convergence of repeated
powers of a single fixed-algebra channel.  Sparse/expander models are known
routes to fast scrambling \cite{BarbonMagan2012,BentsenGuLucas2019}, and OTOC
decay is a standard route from Hamiltonian dynamics to channel decoupling
\cite{HosurQiRobertsYoshida2016,YoshidaKitaev2017}.  What remains for a named
deterministic Hamiltonian family is a bound or scaling analysis showing that
the induced weak-emission channel has the Page entropy curve, the pre- and
post-Page second R\'enyi decoupling diagnostics, and the expected failure of
free, integrable, or resonant controls.

\section{Page and Replica Contractions}
\label{app:page-replica}

This appendix gives the finite-dimensional Page and replica formulas used in
Section~\ref{sec:information-flow}.  For a Haar-random pure state on
$X\otimes Y$, with
$d_X=\dim X$, $d_Y=\dim Y$, and $d_X\le d_Y$, Page's formula is
\begin{equation}
    \langle S_{\rm vN}(X)\rangle
    =
    \sum_{k=d_Y+1}^{d_Xd_Y}\frac1k
    -
    \frac{d_X-1}{2d_Y}
    =
    \log d_X-\frac{d_X}{2d_Y}+O(d_Y^{-2}).
\end{equation}
Approximate unitary two-designs reproduce the second moments needed for
Page-purity and code-subspace decoupling
\cite{DankertCleveEmersonLivine2009,SzehrDupuisTomamichelRenner2013};
fixed R\'enyi-$n$ moments require the corresponding $n$th moments
\cite{HarrowLow2009}.

For a Haar-typical pure state on a $b\times r$ factorization,
\begin{equation}
    \E\,\Tr\rho_R^2
    =
    \frac{b+r}{br+1}
    =
    \frac1r+\frac1b+O((br)^{-1}).
\end{equation}
With
$b=D_{\rm B}(E)=e^{\Smicro(E)}$ and
$r=D_{\rm rad}(E)=e^{\Delta S_{\rm rad}(E)}$, this gives
Eq.~\eqref{eq:renyi-two-island}.  For fixed higher R\'enyi moments, the Haar
moment expansion has the standard permutation form
\begin{equation}
    \E\,\Tr\rho_R^n
    =
    \frac{1}{(br)_n}
    \sum_{\sigma\in S_n}
    b^{C(\sigma)}r^{C(\tau\sigma)}.
\end{equation}
Here $\tau=(12\cdots n)$ is the cyclic permutation imposed by
$\Tr\rho_R^n$, and $C(\sigma)$ is the number of cycles of $\sigma$.  The
identity contraction gives the Hawking/no-island branch, while
$\sigma=\tau^{-1}$ gives the post-Page branch controlled by the remaining core
dimension.  This is the same random-matrix replica mechanism that appears in
Page theorem and random-dynamics discussions of island-like behavior
\cite{Urbach2021,deBoerHollanderRolph2024,BasuWenZhou2025}.

The Hamiltonian assigns Hawking weights to radiation histories.  Let
$\alpha$ label a coarse radiation history, including emitted energies,
channels, and time bins, with probabilities $p_\alpha$ supplied by the
weak-coupling branching rates.  A random purification of this weighted
history distribution gives, at large dimension,
\begin{equation}
    \E\,\Tr\rho_R^n
    \simeq
    \sum_{\sigma\in S_n}
    b^{C(\sigma)-n}
    \prod_{c\in{\rm cycles}(\tau\sigma)}
    \Tr p^{|c|}.
\end{equation}
For $p_\alpha=1/r$ this reduces to the Haar/Page expression.  The two leading
contractions are
\begin{equation}
    \sigma=1:
    \qquad
    \E\,\Tr\rho_R^n\simeq \Tr p^n,
\end{equation}
and
\begin{equation}
    \sigma=\tau^{-1}:
    \qquad
    \E\,\Tr\rho_R^n\simeq b^{1-n}.
\end{equation}
Thus the no-island branch is the R\'enyi entropy of the emitted Hawking
history distribution, while the island branch is the remaining core state
count.  At second R\'enyi order this gives
Eq.~\eqref{eq:weighted-renyi-two-main}.

\section{Finite-Dimensional Check}
\label{app:finite-check}

This appendix gives a small finite-dimensional check of the
information-theoretic step used in the main text.  It checks the conditional
statement that, given a Page-typical emission map and the shrinking
Hilbert-space dimensions, the radiation entropy follows Page's formula.  It
also illustrates that matching a coarse thermal energy-bin distribution is a
weaker condition than producing fine-grained radiation entropy.  It is not a
test of whether a particular sparse Hamiltonian generates the Page-typical
maps dynamically.

The example uses a coarser regularization adapted to the entropy scaling.  The
integer $L$ is an area or entropy label, unrelated to the jump operator
$L_{E;\omega\lambda\mu}$, and $q$ is the Hilbert-space dimension per entropy
unit:
\begin{equation}
    \dim\mathcal H_L=q^{L^2}.
\end{equation}
One regulated emission step factors the current space as
\begin{equation}
    \mathcal H_L\simeq \mathcal H_{L-1}\otimes\mathcal H_{\rm shell}(L),
    \qquad
    \dim\mathcal H_{\rm shell}(L)=q^{2L-1}.
\end{equation}
The removed shell is the outgoing radiation factor for that step.  Its energy
measurement is a coarse graining
\begin{equation}
    \mathcal H_{\rm shell}(L)=\bigoplus_h \mathcal H_h(L),
\end{equation}
where $h$ labels energy bins.  The integer bin dimensions are chosen to
approximate the four-dimensional thermal weights
\begin{equation}
    P_h\propto\int_{x_h}^{x_{h+1}}dx\,x^2e^{-x},
    \qquad x=\beta\omega.
\end{equation}
Thus $\dim\mathcal H_h(L)/\dim\mathcal H_{\rm shell}(L)$ is the prescribed
coarse-grained thermal probability $P_h$ for energy bin $h$, up to integer
rounding.  The full shell Hilbert space carries the fine-grained purification
data.

For the finite-dimensional example we used $q=2$ and $L_0=4$, evolved through
$L=4\to3\to2\to1$, and compared three classes of initial states on
$\mathcal H_{L_0}$.  A Haar-random state is equivalent, for this fixed
factorization, to applying a Haar-random emission isometry to a fixed input.
A basis state is a non-typical state with no radiation entanglement.  A flat
product state has the right energy-bin probabilities and no fine-grained
purification structure.  Five Haar-random samples were used; the two reference
states are deterministic for this factorization.

In Table~\ref{tab:page-check}, $S_{\rm rad}$ means $S_{\rm vN}(R)$ for the
radiation factors emitted so far.  For a Haar-random state, the expected
radiation entropy is the entropy of the smaller factor: the emitted radiation
at early times and the remaining core after the Page crossing.

\begin{table}[H]
\centering
\small
\begin{tabular}{lcccc}
\toprule
step & $L$ transition & smaller factor & calculated $S_{\rm rad}$ & expected $S_{\rm Page}$ \\
\midrule
1 & $4\to3$ & radiation & $4.7269$ & $4.7270$ \\
2 & $3\to2$ & core & $2.7707$ & $2.7706$ \\
3 & $2\to1$ & core & $0.6931$ & $0.6931$ \\
\bottomrule
\end{tabular}
\caption{Expected Page entropy versus calculated radiation entropy for the
Haar-random emission isometry.  Entropies are natural-log entropies and Haar
entries are averages over five samples.  The largest absolute deviation from
Page's formula among the five samples is below $9\times10^{-4}$.}
\label{tab:page-check}
\end{table}

The second comparison shows what coarse radiation thermality leaves
unresolved.  Let $d_{\rm TV}$ denote the mean total-variation distance, over the
three emissions, between the measured energy-bin distribution and the target
distribution with weight $x^2e^{-x}$.  At the final step Page's formula gives
$S_{\rm Page}=0.6931$.  Table~\ref{tab:reference-state-check} compares the
Haar-random state with the two reference states.  The zeros in the last two
rows are exact for these product states: they carry no radiation entanglement
and no mutual information between the previously emitted radiation $R_{<j}$
and the newly emitted radiation factor $R_j$.

\begin{table}[H]
\centering
\small
\begin{tabular}{lcccp{2.05in}}
\toprule
input state
& $S_{\rm rad}$
& $d_{\rm TV}$
& $\max_j I(R_{<j}:R_j)$
& interpretation
\\
\midrule
Haar-random
& $0.6931$
& $0.086$
& $5.414$
& Page curve and post-Page correlations
\\
basis
& $0$
& $0.986$
& $0$
& neither energy-bin thermality nor radiation entanglement
\\
flat product
& $0$
& $0.085$
& $0$
& energy-bin thermality without radiation entanglement
\\
\bottomrule
\end{tabular}
\caption{Comparison at the end of the three-emission sequence.  The flat
product state has energy-bin statistics close to the Haar-random state while
carrying no radiation entropy or mutual information between emitted radiation
factors.}
\label{tab:reference-state-check}
\end{table}

The mutual information $I(R_{<j}:R_j)$ is zero at the first emission and
becomes nonzero after the Page crossing.  The basis state fails both the
spectral and fine-grained comparisons.  The flat product state matches the
coarse energy-bin distribution while carrying no radiation entropy or mutual
information between emitted radiation factors.  Thus coarse energy-bin
thermality and Page-curve behavior probe different parts of the emitted
radiation state.

\bibliographystyle{unsrturl}
\bibliography{refs}

\end{document}
